% Source PDF: source/普通高考/2011/2011新课标理(河南,黑龙江,吉林,海南,宁夏,山西,新疆).pdf, page 1, question 3.
% pdfimages found no matching embedded bitmap; comparison uses a 300dpi page
% render and a 100/50-pixel grid crop at tmp/review_flowchart_2011_new_curriculum_science/.
% Elements: rounded start/end terminals, input/output parallelograms, three
% process rectangles, one decision diamond, and the loop-back connector. All
% borders/connectors are solid; no dashed or hidden segments are present.
% Node -> directed-edge list: start -> input N -> (k=1,p=1) -> (p=p*k)
% -> (k<N); no -> output p -> finish; yes -> (k=k+1) -> loop entry before
% p=p*k. The upward yes arrow and leftward loop-entry arrow are independent.
\begin{tikzpicture}[
  >=Stealth,
  x=1cm,y=1cm,
  line width=0.45pt,
  line cap=round,line join=round,
  font=\normalsize,
  flow/.style={draw,anchor=center,align=center,inner xsep=4pt,inner ysep=2.5pt,
    execute at begin node={\setlength{\parindent}{0pt}}},
  terminal/.style={flow,rounded corners=.18cm,minimum width=1.40cm,minimum height=.74cm},
  process/.style={flow,minimum height=.78cm},
  io/.style={flow,shape=trapezium,trapezium left angle=72,trapezium right angle=108,
    minimum width=2.25cm,minimum height=.80cm},
  decision/.style={flow,diamond,aspect=2.8,minimum width=4.15cm,minimum height=1.05cm,inner sep=0pt},
  branchlabel/.style={inner xsep=1pt,inner ysep=1pt,anchor=center,align=center,
    execute at begin node={\setlength{\parindent}{0pt}}},
]
  % Safety margin keeps the rounded finish and loop line inside the bbox.
  \path[use as bounding box] (-2.25,-8.35) rectangle (4.05,0.42);

  \node[terminal] (start) at (0,0) {开始};
  \node[io,minimum width=2.25cm] (input) at (0,-1.20) {输入 $N$};
  \node[process,minimum width=3.25cm] (init) at (0,-2.48) {$k=1,p=1$};
  \node[process,minimum width=2.10cm] (multiply) at (0,-3.78) {$p=p\cdot k$};
  \node[process,minimum width=2.35cm] (increment) at (2.65,-3.78) {$k=k+1$};
  \node[decision] (test) at (0,-5.12) {$k<N$};
  \node[io,minimum width=2.25cm] (output) at (0,-6.45) {输出 $p$};
  \node[terminal] (finish) at (0,-7.75) {结束};
  \coordinate (loop-merge) at (0,-3.08);

  \draw[->] (start.south) -- (input.north);
  \draw[->] (input.south) -- (init.north);
  \draw (init.south) -- (loop-merge);
  \draw[->] (loop-merge) -- (multiply.north);
  \draw[->] (multiply.south) -- (test.north);
  \draw[->] (test.south) -- node[branchlabel,right] {否} (output.north);
  \draw[->] (output.south) -- (finish.north);

  % True branch rises into k=k+1, then returns left into the p update entry.
  \coordinate (yes-right) at (2.65,-5.12);
  \coordinate (loop-top) at (2.65,-3.08);
  \draw (test.east) -- node[branchlabel,above] {是} (yes-right);
  \draw[->] (yes-right) -- (increment.south);
  \draw (increment.north) -- (loop-top);
  \draw[->] (loop-top) -- (loop-merge);
\end{tikzpicture}
